% Compile from this directory with: xelatex why_three_qiskit7.tex
\documentclass[11pt,a4paper]{article}
\usepackage{kotex}
\usepackage{amsmath,amssymb,booktabs,geometry,hyperref}
\geometry{margin=25mm}
\hypersetup{colorlinks=true,urlcolor=blue}
\title{I01 Qiskit 실험 7: 확률적 위상 브리지와 기준점 사이의 불확실성 바닥}
\author{RELATIONAL ATLAS --- 시간적 위상 기준의 비용 지도}
\date{2026년 10월 4일}

\begin{document}
\maketitle

\begin{abstract}
앞선 결정론적 에일리어싱 실험은 기준점 사이의 드리프트가 숨을 수
있음을 보였다. 이번에는 위상 영점의 시간 변동을 명시적인 확률적
토이 모형으로 바꾸고, 과학 측정 전후의 기준점이 가까워져도 얼마의
불확실성이 남는지 계산했다. 알려진 영평균 지수상관 공분산을 가정하면
완벽하게 측정한 두 기준점 사이에서도 과학 시점의 위상에 양의 조건부
분산이 남는다. 개별 블록의 미지 잔차를 자유 nuisance로 처리한 정확
Fisher rank는 여전히 3이다. 반면 확률 법칙과 블록 독립성을 추가
가정하면 여러 블록을 평균해 공통 과학 위상을 통계적으로 추정할 수
있다. 이때 드리프트 분산을 빼고 샷 오차만으로 만든 95\% 구간은
심하게 과소포괄된다. Aer 회로와 별도의 반복-ensemble 계산으로
구조적 식별성과 확률적 추정 성능을 분리했다.
\end{abstract}

\section{기준점 사이의 확률적 위상 브리지}
과학 측정은 시간 $t=0$, 알려진 위상의 기준 측정은 $t=-h,+h$에
놓는다. 위상 드리프트 $\delta(t)$는 영평균, 분산 $\sigma^2$,
공분산
\begin{equation}
 \operatorname{Cov}[\delta(t),\delta(s)]
 =\sigma^2\exp(-|t-s|/\ell)
 \label{eq:cov}
\end{equation}
인 가우스 과정이라고 \emph{가정}한다. 이는 실제 장치의 잡음 법칙을
측정하거나 입증한 결과가 아니라, 상관시간 $\ell$에 대한 명시적
반례/비용 모형이다. $r=e^{-h/\ell}$로 두면 두 기준점의 상관은
$r^2$이고, 가우스 조건부 평균과 분산은
\begin{align}
 \mathbb E[\delta(0)\mid\delta(-h),\delta(+h)]
 &=\frac{r}{1+r^2}\,[\delta(-h)+\delta(+h)],\\
 v_{\rm bridge}
 &=\operatorname{Var}[\delta(0)\mid\delta(-h),\delta(+h)]
 =\sigma^2\frac{1-r^2}{1+r^2}
 =\sigma^2\tanh(h/\ell).
 \label{eq:bridge}
\end{align}
따라서 $h>0$이면 기준점 두 곳의 위상을 \emph{무한 shots로} 알아도
중간 위상에는 불확실성이 남는다. 이 값은 기준 측정 자체의 샷
오차를 무시한 낙관적인 하한이다. 본 실험은 $\sigma=0.120$ rad,
$h/\ell=0.05$와 $1.0$을 사전에 정했다.

\section{rank 3과 통계적 추정은 다른 질문}
한 블록의 과학 확률은 $\alpha+\delta(0)$에 의존한다. 기준점
$\delta(\pm h)$는 별도로 측정할 수 있어도 $\delta(0)$의
\emph{실현값}을 자유 nuisance로 허용하면
$(\alpha,\delta(0))\mapsto(\alpha+c,\delta(0)-c)$가 정확한
불변 방향이다. 따라서 비제약 네 축의 nuisance-profiled Fisher
rank는 $\tau=0,\pi/2$ 모두 3이다. 여러 블록을 더해도 각 블록에
자유로운 미지 잔차를 허용하면 같은 null 방향이 남는다.

식~\eqref{eq:cov}의 영평균 확률 법칙과 블록 간 독립성을 알고
신뢰한다면 다른 추론 문제를 만들 수 있다. 블록마다 조건부 평균을
뺀 과학 위상 추정의 오차는 브리지 잔차와 유한-shots 측정 오차의
합이다. 두 성분의 분산을 각각 $v_{\rm bridge},v_{\rm shot}$으로
두면 $B$개 독립 블록 평균의 근사 표준편차는
\begin{equation}
 \operatorname{SD}(\hat\alpha_B)
 \simeq\sqrt{(v_{\rm bridge}+v_{\rm shot})/B}.
 \label{eq:budget}
\end{equation}
이는 확률 \emph{가정}으로 얻은 평균 추정 가능성이지, 한 블록의
정확 Fisher rank가 4로 변했다는 뜻이 아니다. 과학 시점 잔차에
미지의 비영 평균을 별도로 허용하면 그 평균과 $\alpha$는 다시
혼동된다.

\section{정확 브리지 분산과 반복-ensemble 비용}
과학 측정에는 기존 5준위 에너지 매칭 시계, 9개 위상 표지,
$n=4$ 비제약 토이 상태를 쓴다. 각 위상 실행의 조건부 복귀
확률 미분과 시계 후선택률 $1/5$로 유한-shots Fisher 오차를
계산했다. 아래 반복-ensemble은 \emph{Aer 회로 반복이 아니다}.
정확 가우스 브리지와 국소 Fisher의 가우스 샷 오차를 합쳐 조건당
20,000개 독립 ensemble을 생성하고, 각각 $B=8$ 블록 평균의
구간 포괄률을 평가했다. 목표 반폭 $0.02$ rad에 필요한 블록 수도
식~\eqref{eq:budget}으로 계산했다.
\begin{center}
\begin{tabular}{@{}ccccccc@{}}
\toprule
$\tau/\pi$ & $h/\ell$ & 브리지 SD & 샷 SD & 샷만의 포괄률
& 브리지 포함 포괄률 & 목표 블록 수 \\
\midrule
0 & 0.05 & 0.0268 & 0.0071 & 0.380 & 0.948 & 8 \\
0 & 1.00 & 0.1047 & 0.0065 & 0.095 & 0.950 & 106 \\
1/2 & 0.05 & 0.0268 & 0.0262 & 0.827 & 0.949 & 14 \\
1/2 & 1.00 & 0.1047 & 0.0236 & 0.335 & 0.950 & 111 \\
\bottomrule
\end{tabular}
\end{center}
표의 SD는 블록당 rad 단위이고, 포괄률은 $B=8$의 근사 95\%
구간에 대한 값이다. $\tau=\pi/2$의 짧은 간격에서 ``샷만''
포괄률이 높아 보이는 이유는 시계 조건부 측정 자체가 약해 샷
구간이 넓기 때문이다. 정밀도가 더 좋다는 뜻이 아니다. $h/\ell=1$
에서는 드리프트가 샷 오차보다 커서 기준점 shots만 늘려도
오차 바닥을 없앨 수 없다. 이 비용 지도는 $\sigma,\ell$과
영평균 법칙을 \emph{정확히 안다는} 낙관적 계산이다.

\section{Aer 회로의 8블록 대조}
각 $h/\ell$에 독립 브리지 8블록을 사전 seed로 생성했다. 한
블록은 앞 기준, 과학, 뒤 기준의 세 위상 실행이며 각 실행은
검출 축 4개와 표지 9개, 즉 36회로다. $\tau=0$에서 총
$2\times8\times3\times36=1728$개의 Aer statevector 회로를
각각 4096 shots, 최대 CPU 병렬 6으로 실행했다. 각 세그먼트의
위상을 표지별 raw counts로 적합하고 식~\eqref{eq:bridge}의
조건부 평균을 뺐다.
\begin{center}
\begin{tabular}{@{}cccc@{}}
\toprule
$h/\ell$ & 8블록 평균 $\hat\alpha$ & 샷만의 근사 95\% 구간
& 브리지 포함 근사 95\% 구간 \\
\midrule
0.05 & 0.2954 & $[0.2905,0.3004]$ & $[0.2762,0.3147]$ \\
1.00 & 0.3745 & $[0.3702,0.3788]$ & $[0.3018,0.4472]$ \\
\bottomrule
\end{tabular}
\end{center}
참값은 $\alpha=0.3000$ rad다. 넓은 간격의 선택된 단일 8블록
표본은 브리지 포함 구간에서도 참값을 약 $0.0018$ rad 차이로
놓쳤다. 사전 seed를 결과에 맞춰 바꾸지 않았다. 반복-ensemble의
약 95\% 포괄률은 모든 개별 구간이 참값을 포함한다는 약속이
아니다. 반면 샷만의 구간은 이 표본에서 현저히 과신한다.

\section{$P=0$ 유한표본 대조군의 별도 품질 점검}
실험 6의 $P=0$ Aer 대조군 16 seed$\times$2 시계 위상에서
32개 명목 5\% 재배분 검정 중 5개가 양성이었다. 정확한 $P=0$
확률은 공통 위상 변화에 불변이다. 원자료의 두 arm 차이를 각
회로의 정확 이항분산으로 표준화하니 평균 $-0.0033$, 표본분산
$1.0683$이었다. 정확 이항모형에서 만든 5000개 동형 패널의
분산 95\% 범위는 $[0.9193,1.0843]$이고, 관측 분산의 상측 꼬리
비율은 $0.0526$이었다. 평균 차이의 양측 꼬리 비율은 $0.916$이다.
즉 명확한 방향성 편향은 보이지 않지만 변동성은 이 점검의
상단에 가깝다. 32개 검정의 양성 수 자체는 독립 5\% 가정 아래
상측 확률 약 $0.020$으로 다소 많다. 이 둘만으로 Aer 결함 또는
물리 효과를 확정하지 않으며, 대조군 보정은 여전히 점검 사항이다.

\section{판정과 다음 게이트}
기준점 간격을 줄이면 확률적 드리프트의 불확실성은 줄지만,
과학/기준 측정 시점이 분리된 한 한 블록의 미지 실현값은 정확히
결정되지 않는다. 통계 법칙을 이용해 여러 블록을 평균할 수
있지만 그 가정과 분산을 오차 예산에 포함해야 한다. 본 실험의
가우스 지수상관, 블록 독립성, 완벽한 시계 및 기준 위상은 모두
가정이다. 이 결과 역시 자발적 세 공간 차원 선택의 증거가 아니다.
후속으로는 상관시간과 평균의 미지성, 비가우스 드리프트 및 기준
측정 자체의 위상 오차를 분리해 검토해야 한다. QPU는 사용하지 않았다.

\paragraph{재현성.} 회로/브리지 코드
\texttt{../src/relational\_atlas\_i01/phase\_reference/qiskit\_clock\_stochastic\_bridge.py}, 회귀 검사
\texttt{../tests/test\_qiskit\_clock\_stochastic\_bridge.py}, 원자료
\texttt{../results/qiskit\_clock\_stochastic\_bridge.json}.
$P=0$ 진단 코드/원자료는
\texttt{../src/relational\_atlas\_i01/phase\_reference/qiskit\_pzero\_sampling\_audit.py}와
\texttt{../results/qiskit\_pzero\_sampling\_audit.json}이다.
시작 seed 20261010, Qiskit 2.4.2, Qiskit Aer 0.17.2.
Windows Aer의 합성 \texttt{multiplexer} 접근 위반을 피하려고
준비 회로를 $u3,cx$로 분해했다. TeX 엔진이 없어 PDF 렌더링은
아직 검증하지 못했다.
\end{document}
